% !TeX root = ../../tesi.tex

\section{Metodo di lavoro}

\subsection{Pianificazione e interazione con il \emph{tutor} aziendale}

Confronto le modalità di interazione previste dal piano di lavoro con quelle effettive: le riunioni
quotidiane, i resoconti di fine giornata su Slack, la riunione settimanale, i confronti su richiesta
e le giornate in sede. Discuto le difficoltà della fase iniziale e il modo in cui le ho superate.
Descrivo infine i resoconti settimanali al relatore e la svolta della terza settimana: con gli
obiettivi obbligatori già soddisfatti, io e il \emph{tutor} abbiamo deciso di dedicare il resto
dello \emph{stage} alla sperimentazione. Un diagramma temporale mette a confronto la pianificazione
originaria del piano di lavoro con il modo in cui ho usato davvero il tempo, attività per attività.

\subsection{Organizzazione del lavoro e tracciamento delle attività}

Illustro come ho organizzato il mio lavoro, svolto in gran parte da remoto: come individuavo le
attività e ne stabilivo le priorità, come ne tenevo traccia e come preparavo le modifiche più estese
prima di realizzarle.

\subsection{Strumenti e assistenti di programmazione}

Presento gli strumenti che ho impiegato per lo sviluppo, la verifica e il rilascio, tra cui gli
assistenti di programmazione basati su modelli linguistici. Di questi descrivo l’uso che ne ho
fatto: la sperimentazione e la validazione rapida delle idee, fino a prototipi funzionanti nell’arco
di una giornata, a cui seguiva l’implementazione definitiva scritta a mano; la stesura dei documenti
di progettazione, per individuare in anticipo i casi limite e i \emph{test} necessari; le attività
di supporto, come la generazione dei dati di prova. Ciò che ho imparato da questo uso lo discuto
nella \secref{sec:maturazione}.
